\chapter{Queries}
\label{ch:queries}

A rule reacts; a query asks. \fn{defquery} defines a parameterized
pattern set that is run on demand against working memory and returns the
facts that match, and \fn{defgraphquery} does the same against a graph of
nodes and edges loaded from a file. Both use the left-hand-side syntax of
rules, compiled into their own part of the network, so a query costs
nothing until it is run.

\section{defquery}

\begin{lstlisting}
(defquery <name> ["comment"]
  (declare (variables ?v1 ?v2 ...))
  <conditional element>*)
\end{lstlisting}

The \fn{declare} block is mandatory and names the query's parameters.
The patterns use them like any variable, and a parameter may appear in
several patterns; the query returns every combination of facts that
matches when the parameters take the values passed to \fn{run-query}.
From \file{samples/defquery/queries.clp}:

\begin{lstlisting}
(defquery queryByTransactionStart (declare (variables ?trxId ?startTime))
    (ApplicationRequestLog
        (TransactionID ?trxId)
        (Timestamp ?startTime)))

(defquery queryByTransactionAfterStart (declare (variables ?trxId ?startTime))
    (ApplicationRequestLog
        (TransactionID ?trxId)
        (Timestamp ?ts&:(> ?ts ?startTime))))
\end{lstlisting}

The first query finds the log entry of a transaction at an exact time;
the second every entry of the transaction after a time, using the
parameter inside a predicate. \file{samples/defquery/run.clp} loads the
templates and queries, loads a data file and runs them:

\begin{lstlisting}
(batch ./samples/defquery/templates.clp)
(batch ./samples/defquery/queries.clp)
(load-facts ./samples/defquery/log_data.dat)
(bind ?result (run-query queryByTransactionStart "1234baaa" 1269285900120))
(printout t ?result crlf)
\end{lstlisting}

\begin{lstlisting}[style=shell]
f-3 (ApplicationRequestLog (IPAddress "129.0.0.1") (Hostname "localhost")
    (ServerType "Tomcat") (PortNumber 8080) (Timestamp 1269285900120) ...
    (TransactionID "1234baaa") (ThreadID "2") )
\end{lstlisting}

\fn{run-query} takes the query name and one value per declared variable,
in order, and returns a list with one entry per match. Each entry is the
tuple of facts that matched, one fact per pattern; the shell prints a
tuple per line, and \fn{printout} prints the facts one after another.
Bind the list and pass it to the list functions (\fn{length\$},
\fn{nth\$}), or print it. \fn{(query-time name)} reports how long the
last run took in milliseconds, and \fn{(watch-query name)} traces a
query's evaluation.

Patterns join through shared variables exactly as in a rule, and a
predicate may relate a pattern to an earlier one:

\begin{lstlisting}[style=shell]
Morendo> (clear)
true
Morendo> (deftemplate guest (slot name) (slot hobby))
true
Morendo> (assert (guest (name "a") (hobby "x")))
2
Morendo> (assert (guest (name "b") (hobby "x")))
3
Morendo> (defquery same-hobby (declare (variables ?h))
  (guest (name ?n1) (hobby ?h))
  (guest (name ?n2&:(neq ?n2 ?n1)) (hobby ?h)))
true
Morendo> (run-query same-hobby "x")
f-3 (guest (name "b") (hobby "x") ) f-2 (guest (name "a") (hobby "x") )
f-2 (guest (name "a") (hobby "x") ) f-3 (guest (name "b") (hobby "x") )
\end{lstlisting}

A query's facts are those in working memory at the time it runs; unlike
a rule it keeps no activation, and its network is emptied before every
run, so running it again after facts change gives the new answer.
Queries and rules do not share alpha nodes, so a query never slows down
rule matching.

\section{Graph queries}

\Morendo\ declares three templates at start-up, \fn{Node}, \fn{Edge} and
\fn{Graph}, with the slots of the classes in \fn{org.morendo.model}: a
node has \fn{id}, \fn{label}, \fn{nodeType}, \fn{reference} and
\fn{weight}; an edge has \fn{id}, \fn{label}, \fn{source}, \fn{target}
and \fn{weight}. A graph file is a list of node and edge facts,
\file{samples/graphquery/graphdata.clp}:

\begin{lstlisting}
(Node (id "1")(label "electric vehicle")(nodeType "concept"))
(Node (id "2")(label "autonomous vehicle")(nodeType "concept"))
(Node (id "3")(label "tesla model s")(nodeType "noun"))
(Node (id "4")(label "leaf")(nodeType "noun"))
(Edge (id "100")(source "3")(target "1"))
(Edge (id "102")(source "4")(target "1"))
(Edge (id "105")(source "3")(target "2"))
\end{lstlisting}

\fn{(load-graph file)} reads such a file and returns the facts as a list
\emph{without} asserting them into working memory. A graph query is a
\fn{defquery} over \fn{Node} and \fn{Edge} patterns, joined through the
id slots; \file{samples/graphquery/query.clp}:

\begin{lstlisting}
(defgraphquery conceptDistanceOne (declare (variables ?label))
  (Node (nodeType "noun") (label ?label) (id ?sourceid))
  (Edge (source ?sourceid) (target ?targetid))
  (Node (nodeType "concept") (id ?targetid)))
\end{lstlisting}

It reads: starting from the noun with the given label, follow one edge
and return the concept nodes reached. \fn{run-graph-query} takes the
graph, the query name and the parameter values:

\begin{lstlisting}
(batch ./samples/graphquery/query.clp)
(bind ?data1 (load-graph ./samples/graphquery/graphdata.clp))
(bind ?query1 (run-graph-query ?data1 conceptDistanceOne "tesla model s"))
(printout t "first query results: " ?query1 crlf)
\end{lstlisting}

\begin{lstlisting}[style=shell]
first query results: f-2 (Node (id "1") (label "electric vehicle") (nodeType "concept")
  (reference "nil") (weight "nil") ) f-3 (Node (id "2") (label "autonomous vehicle")
  (nodeType "concept") (reference "nil") (weight "nil") )
\end{lstlisting}

Longer paths are longer chains of \fn{Edge} and \fn{Node} patterns;
\fn{conceptDistanceTwo} in the same file follows two edges to reach
project nodes. Because the graph is held outside working memory, several
graphs can be loaded and queried independently in one engine, and rules
are unaffected by them. To reason over a graph with rules instead, load
it with \fn{load-facts}: the same file asserts ordinary \fn{Node} and
\fn{Edge} facts.
